\writeSectionFrame{\presentationTitle}{Fazit}
\begin{frame}{Anwendungsfälle}
	\begin{itemize}
 		\item Regel-Engines
 		\item Formel-Generatoren
 		\item Daten-Generatoren
 		\item Domain Specific Languages
 		\item Filter
 	\end{itemize}
\end{frame}

\note{
	Regeln sind im allgemeinen sehr schwierig, weil sie die gesamte Bandbreite an Operatoren, Variablen und allen möglichen
	Ausdrücken anbeiten. Sie können jedoch nicht fest codiert werden, weil sie flexibel angepasst werden können. Solche
	Regelwerke können auch zur Beschreibung von Geschäftsprozessen verwendet werden. Formel-Generatoren funktionieren
	ähnlich wie Regel-Engines, nur dass im Allgemeinen zu Zahlen ausgewertet wird. Wenn Daten mithilfe von Definitionen
	beschrieben werden können, kann dieses Muster auch zur Generierung von Daten verwendet werden. Domain-spezifische Sprachen
	sind Sprachen für eine spezielle Aufgabenstellung. Das Interpreter-Muster ermöglicht es, dass deren Grammatik
	beliebig formuliert werden kann. Filter-Ausdrücke sollen in beliebigen Kombinationen formuliert werden können. Auch hier 
	kann das Interpretermuster eingesetzt werden. 
}

\vorteil{Erweiterbarkeit}
\note{
    Es ist einfach, die Grammatik durch das Hinzufügen neuer Regeln zu erweitern. Neue Expression-Unterklassen können problemlos hinzugefügt werden, ohne die bestehenden Klassen zu verändern.
}

\vorteil{Modularität}
\note{
    Jeder Ausdruck in der Grammatik wird durch eine eigene Klasse dargestellt, was das System modular und leichter verständlich macht.
}

\vorteil{Leichte Implementierung für kleine Sprachen}
\note{
    Für relativ einfache und kleine Sprachen oder Grammatikdefinitionen ist das Muster einfach zu implementieren.
}

\vorteil{Single-Responsibility-Prinzip}
\note{
    Jede Ausdrucksklasse kümmert sich nur um ihren Teil des Gesamtausdrucks. 
}

\nachteil{Leistungseinbußen}
\note{
    Da jede Regel durch eine eigene Klasse repräsentiert wird, können viele kleine Objekte erzeugt werden. Für komplexe Sprachen oder große Eingaben kann das zu einem erheblichen Leistungsverlust führen.
}

\nachteil{Komplexität für größere Sprachen}
\note{
    Das Muster ist nicht gut geeignet für komplexe oder umfangreiche Sprachen. Die Anzahl der Klassen kann explodieren, was zu einer unübersichtlichen und schwer wartbaren Codebasis führt.
}

\nachteil{Schwierige Optimierung}
\note{
    Es ist schwer, eine optimierte Ausführung oder Optimierungen wie Just-in-Time-Compilation (JIT) oder andere fortgeschrittene Techniken in ein Interpreter-Muster zu integrieren, denn die Optimierungstechniken sind schwer auf viele Klassen umzusetzen. Auch Techniken wie Vor-Kompilierung sind schwierig umzusetzen, wenn die gesamte Interpretation zur Laufzeit stattfindet. 
}

\nachteil{Nicht für jede Anwendung geeignet}
\note{
    Das Muster ist ideal für einfache formale Sprachen oder spezielle Probleme, aber für generelle Programmiersprachen oder sehr komplexe Interpretationsaufgaben könnte ein dedizierter Parser/Compiler besser geeignet sein.
}

\nachteil{Viele Klassen bei komplexen Grammatiken}
\note{
	Komplexe Grammatiken können schwer zu warten werden, weil mindestens eine Klasse für jede Regel in der Grammatik definiert
	wird.
}